\section{多元函数微分学}
\subsection{多元函数的概念、极限与连续性}
\begin{mydef}{二元函数}{two-variable function}
    设\(D\)是平面上的一个点集，如果对每个点\(P(x, y) \in D\)，都有一个唯一的数\(z = f(x, y)\)，则称\(f(x, y)\)是\(D\)上的一个二元函数。
    \(D\)为该函数的定义域，\(z = f(x, y) \in D\)是该函数的值域。
\end{mydef}
\begin{conclusion}{二元函数的几何意义}{two-variable function geometry}
    空间点集\(\{(x, y, z) | z = f(x,y) \in D\}\)为二元函数\(z = f(x, y)\)的图形。通常它是一张曲面。
    曲面\(z = f(x, y)\)与平面\(z = c\)的交线为\(z = f(x, y) = c\)，这是一个曲线。
    曲线\(z = f(x, y) = c\)在\(Oxy\)平面上的投影曲线:\(f(x, y) = C\)称为\(z = f(x, y)\)的\textbf{等高线}
\end{conclusion}
\begin{conclusion}{一元函数与二元函数的联系与区别}{one-variable function and two-variable function}
    \begin{enumerate}
        \item 一元函数是二元函数的特殊情形：让一自变量动，另一自变量固定，或让\((x, y)\)沿某曲线变动，二元函数就转换为一元函数
        \item 一元函数中，自变量\(x\)代表直线上的点，只有两个变动方向，而二元函数中，自变量\(x, y\)代表平面上的点，有无数个变动方向。
        \item 一元函数\(z = f(x)(a < x <b)\)也可看成二元函数，其定义域为\(
            a < x < b, -\infty < y +\infty
        \)
    \end{enumerate}
\end{conclusion}
\begin{mydef}{二元函数极限}{two-variable function limit}
    设函数\(f(x, y)\)在开区域或闭区域\(D\)内有定义，\(M_0(x_0, y_0)\)是\(D\)的内点或边界点
    则\\
    \(
        \lim\limits_{\substack{x \to x_0 \\ y \to y_0}} = 
        \lim\limits_{(x, y) \to (x_0, y_0)}
            f(x, y) = A  \Leftrightarrow
        \forall \varepsilon > 0, \exists \delta > 0,
        \text{当}f(x, y) \in D, \text{且}0 < \sqrt{(x - x_0)^2 + (y - y_0)^2} < \delta
        \text{时，有}\left| f(x, y) - A \right| < \varepsilon
    \)\\
    这里的极限过程是点\(x, y\)在\(D\)内沿任何路径以任何方式趋于点\((x_0, y_0)\)
\end{mydef}
\begin{conclusion}{二元函数极限与一元函数极限比较}{two-variable function limit and one-variable function limit}
    \begin{enumerate}
        \item 二元函数与一元函数有相同的极限运算法则与极限性质
        \\
        求二元函数极限的常用方法有
        \begin{enumerate}
            \item 用极限的运算法则
            \item 放大缩小法，变量变换法转为求简单的极限或一元函数的极限
        \end{enumerate}
        \item 二元函数\(z=f(x, y)\)与一元函数之极限存在性问题的区别
        \\
        与一元函数极限中"函数在某一点的极限当且仅当它在该点的左右极限存在且相等"相对应。\\
        在二元函数的极限中，极限 \(\lim\limits_{(x, y) \to (x_0, y_0)}
            f(x, y) = A\)存在的充要条件是当点\(P(x, y)\)在定义域内沿任何路径以任何方式趋于点\(P(x_0, y_0)\)时，均有\(
                \lim\limits_{(x, y) \to (x_0, y_0)}
            \)存在且等于\(A\)。 \\
            \\
            若在定义域内沿两条不同路径\((\text{如},\text{直线}y - kx, \text{抛物线}y = x^2)\),极限\(
                \lim\limits_{(x, y) \to (x_0, y_0)} f(x, y)
            \)的值不相等或沿某一路径极限\(\lim\limits_{(x,y) \to (x_0, y_0)} f(x, y)\)不存在，则\(
                \lim\limits_{(x, y) \to (x_0, y_0)} f(x, y)
            \)不存在。这是证明二元函数极限不存在的一个有效方法
    \end{enumerate}
\end{conclusion}
\begin{mydef}{二元函数连续性}{two-variable function continuity}
    设\(z = f(x, y)\)定义在区域\(D\)上，\(P_0(x_0, y_0)\)是\(D\)的内点或边界点，若\(
        \lim\limits_{(x, y) \to (x_0, y_0)} f(x, y) = f(x_0, y_0)
    \)，则称\(f(x, y)\)在\(P_0(x_0, y_0)\)处连续。若\(\forall (x, y) \in D, \text{且}
        \lim\limits_{(x, y) \to (x_0, y_0)} f(x, y) = f(x_0, y_0)
    \)，则称\(f(x, y)\)在\(D\)上连续,即\(f(x,y)\)在\(D\)的所有点上都连续。
\end{mydef}
\begin{conclusion}{二元连续函数的性质}{two-variable function continuity property}
    由自变量\(x\)的初等函数及自变量\(y\)的初等函数经过有限次四则运算和复合运算得到的二元函数
        \(z = f(x, y)\)称为二元初等函数，类似地可定义多元初等函数。\textbf{多元初等函数在其定义区域上是连续的} \\
        判断二元函数连续性有和一元函数相同的方法 \\
        二元连续函数具有以下相应性质
        \begin{enumerate}
            \item 设\(z = f(x,y )\)在\(M(x_0, y_0)\)连续，\(f(x_0, y_0) > 0(< 0)\),则
            \(\exists \delta > 0, \text{当}0 < \sqrt{(x - x_0)^2 + (y - y_0)^2} < \delta
            \text{且} (x, y) \in D \)时，\(f(x, y) > 0(< 0)\)
            \item 最大值最小值定理 \\
            设\(z = f(x, y)\)在有界闭区域\(D\)上连续，则它在\(D\)上一定有最大值和最小值，
            即\(\exists M_1,M_2 \in D,\)使得\(\forall M \in D\),有
            \[
                f(M_2) \leq f(M) \leq f(M_1)
            \]
            其中\(f(M_2)\)为最小值，\(f(M_1)\)为最大值
            \item 中间值定理 \\
            设\(z = f(x, y)\)在连通的区域\(D\)上连续，\(\forall M_1, M_2 \in D\) ，
            \(
                f(M_1) < f(M_2)
            \),则对任意的函数\(\mu\),\(f(M_1) < \mu < f(M_2) \)，
            则存在一点\(P(x, y)\)，使得\(f(P) = \mu\)
        \end{enumerate}
\end{conclusion}
\subsection{多元函数的偏导数与全微分}
\begin{mydef}{偏导数}{partial derivative}
\end{mydef}
\begin{conclusion}{偏导数的几何意义}{partial derivative geometric meaning}
\end{conclusion}
\begin{conclusion}{偏导数的计算}{partial derivative calculation}
\end{conclusion}
\begin{mydef}{可微性与全微分}{differentiability and total differential}
\end{mydef}

\begin{conclusion}{\(z = f(x, y)\)在点\((x_0, y_0)\)可微分的必要与充分条件}{all derivative requirement}
\end{conclusion}

\begin{conclusion}{偏导数的连续性，函数可微性，可偏导性与函数连续性的关系}{ inter relationship}
\end{conclusion}

\begin{conclusion}{高阶偏导数、混合偏导数与求导次序无关}{higher order partial derivative and mixed partial derivative}
\end{conclusion}

\begin{conclusion}{多元函数为常数的条件}{constant function}
\end{conclusion}

\begin{conclusion}{多元偏导的求法}{multi-variable partial derivative}

\end{conclusion}

\begin{formula}{全微分四则运算法则}{total differential formula}
\end{formula}
\begin{conclusion}{多元函数的微分法则}{multi-variable function derivative rule}
\end{conclusion}

\begin{conclusion}{复合函数求导法则的注意事项}{compound function derivative rule notice}
\end{conclusion}
\begin{conclusion}{复合函数的二阶偏导数}{second order partial derivative of compound function}
\end{conclusion}

\begin{conclusion}{由一个方程式确定的一元隐函数求导法}{implicit function derivative method}
\end{conclusion}
\begin{conclusion}{由一个方程式确定的二元隐函数求导法}{implicit function derivative method}
\end{conclusion}

\begin{conclusion}{由方程组确定的一元隐函数求导法}{implicit function derivative method}
\end{conclusion}
\begin{conclusion}{由方程组确定的二元隐函数求导法}{implicit function derivative method}
\end{conclusion}
\begin{conclusion}{复合函数求导法则的其他应用}{compound function derivative rule other application}
\end{conclusion}
\subsection{多元函数的极值与最值问题}
\begin{mydef}{多元函数极值与驻点}{multi-variable function extremum and stationary point}
\end{mydef}
\begin{conclusion}{多元函数取得极值的必要条件与充分条件}{multi-variable function extremum necessary and sufficient condition}
\end{conclusion}
\begin{conclusion}{求二元函数极值点的步骤}{two-variable function extremum point}
\end{conclusion}
\begin{conclusion}{条件极值点的必要条件}{condition extremum point necessary condition}
\end{conclusion}

\begin{conclusion}{极值最值问题的提法}{extremum and extremum point}
\end{conclusion}

\begin{conclusion}{求二元函数与三元函数的简单极值问题}{simple extremum problem}
\end{conclusion}

\begin{conclusion}{求二元函数与三元函数的条件极值问题}{condition extremum problem}
\end{conclusion}
\subsection{梯度与方向导数}
\begin{mydef}{方向导数}{directional derivative}
\end{mydef}
\begin{conclusion}{方向导数的存在性与计算公式}{directional derivative existence and formula}
\end{conclusion}
\begin{conclusion}{梯度(向量)与方向导数的最大值}{gradient and directional derivative maximum}
\end{conclusion}
\subsection{多元函数的几何应用}
\begin{conclusion}{用隐式方程表示的曲面}{implicit equation surface}
\end{conclusion}
\begin{conclusion}{用显式方程表示的曲面}{explicit equation surface}
\end{conclusion}
\begin{conclusion}{用参数方程表示的空间曲线}{parameter equation space curve}
\end{conclusion}
\begin{conclusion}{作为两曲线交线的空间曲线}{intersection of two curves}
\end{conclusion}
\subsection{普适方法与技巧}
\begin{conclusion}{讨论二元函数\(z = f(x,y)\)可微的常用方法}{common method of differentiability}
\end{conclusion}